\chapter{Solving It: Splines, Lanczos, and Partitions}
\label{chap:splines-lanczos-partitions}

The previous chapter reached the Euler--Lagrange equation but left it as a condition, not a method: it
said what a stationary path satisfies, and granted, through a single oracle, that the question of
stationarity has an answer. This chapter does the constructive work the equation invites --- it
actually builds the stationary path, and reads off that the path solves the equation. It is the
largest chapter in the book because solving is harder than stating, and it proceeds in four moves: the
cubic spline that suffices, the Lanczos iteration that computes it, the real partitions that scale it
to genuine problems, and the interior node where the second variation first shows its face. One line
from the last chapter holds throughout, and it is worth fixing before we begin: the constructions here
are all built and checked; the one thing that remains granted is the residue-collapse, the same single
oracle, and nothing in the machinery of this chapter quietly adds a second assumption.

\section{The cubic spline suffices}
\label{se:the-cubic-spline-suffices}

The central fact is a law the construction proves and names: cubic-spline sufficiency. Take a path
that is closed-flat in the sense of the ninth chapter --- flat in its slip and admitting it. From such
a path the construction builds a cubic spline, the smooth piecewise-cubic completion through the path's
anchor points, and proves that this spline is \emph{sufficient}: the residue it leaves is exactly zero.
The chain is short and entirely demonstrated. Closed-flatness gives stationarity; stationarity gives
the Euler--Lagrange equation, by the equivalence of the previous chapter; and a path that satisfies the
Euler--Lagrange equation yields a cubic-spline certificate whose residue is the bare truth that true
equals true. To solve the variational problem, in other words, is to draw the cubic spline through the
anchors, and the law of spline sufficiency guarantees that the spline so drawn is a genuine solution.

There is a reason it is the cubic spline specifically, and the experiment on precision states it
exactly. The Law of Spline Sufficiency, as the experiment names it, is that any finite sequence of
anchor events admits a \emph{unique} minimal-curvature completion --- the cubic spline is the one curve
through the anchors that bends as little as possible, and minimal bending is precisely what a
stationary path, with its vanishing first variation, demands. This is also where the chapter's first
honest seam shows. The spline is determinate: once the anchors are fixed, every value between them is
fixed too, with full precision. Whether those values are accurate --- whether the next measured event
falls on the spline --- is a separate question the completion cannot answer, exactly the
precision-against-accuracy distinction met before. The construction supplies the unique sufficient
spline; the world supplies the verdict on whether it was right.

It is worth walking the chain once, because its shape recurs. Suppose the construction is handed a path
it has certified closed-flat. By the ninth chapter that path is stationary; by the tenth it therefore
satisfies the Euler--Lagrange equation; and the construction now builds, from that path, a cubic-spline
certificate --- a record that the piecewise-cubic curve through the anchors meets the equation at every
joint. The law of spline sufficiency reads the residue off that certificate and finds it zero. Notice
what each link costs. The first two are the proved equivalences of the earlier chapters; the
certificate is a finite construction; and the only thing consulted along the way that the construction
could not compute for itself is the oracle's verdict on stationarity, made once, back at the equation.
The spline is the constructive output; the residue of zero is the receipt that it solved the problem;
and the debt, as always, is the single granted line and nothing more.

\section{Lanczos as ordered consumption}
\label{se:lanczos-as-ordered-consumption}

A spline is an answer; computing it is a procedure, and the construction's procedure is the Lanczos
iteration, which it reads as the ordered consumption of knots. The spline's interior structure is a
finite collection of knots --- the joints where its cubic pieces meet --- and the iteration consumes
them one at a time, each step solving for one knot and removing it from the collection still to be
done. Here the construction proves something it genuinely can: the iteration terminates. Each step
strictly shrinks the multiset of remaining knots, a strictly decreasing count over a finite collection
cannot decrease forever, and so the recursion is well-founded --- it reaches the empty collection in a
bounded number of steps and stops. The construction tanges one knot per step, funges the consumed
knots into a growing trace, and the trace provably ends.

What the construction does \emph{not} prove here is the more tempting claim, and the discipline of the
last chapter requires saying so plainly. One would like to assert that the iteration reaches the
\emph{same} terminal residue no matter what order the knots are consumed in --- that the answer does
not depend on the bookkeeping. That order-independence is not a fresh theorem. It is the
residue-collapse restated: the claim that the residual, surveyed across all the directions the knot
ordering exposes, collapses to one value, which is exactly the universal-over-infinite-directions the
oracle was introduced to decide. Termination is built; order-independence of the terminal residue is
granted, and it is granted by the same single oracle and no other. The experiment on the limitation of
indexing marks the boundary from the other side: a recursive catalog may be continued indefinitely and
yet add nothing, returning structurally identical copies of the same content rather than new
information. The Lanczos recursion is the disciplined opposite --- it terminates rather than continuing
emptily --- but the lesson is shared: more iteration is not more truth, and the construction takes from
the iteration only what it has earned, the terminating computation, while marking the one thing it
takes on trust.

The termination argument is worth making concrete, because it is the kind of thing easy to wave at and
important to actually have. Begin with a finite bag of knots --- say five interior joints to be solved.
The first step of the iteration solves one of them and removes it, leaving four; the next leaves three,
then two, then one, then none. The count goes five, four, three, two, one, zero, and a sequence of
counts that strictly decreases and is bounded below by zero cannot run on forever --- it must reach the
floor. That is the whole of the termination proof, and it is genuinely a proof: no appeal to anything
infinite, no oracle, no assumption that the process is well-behaved. A finite bag emptied one item at a
time empties. What the proof does not touch, and cannot, is whether emptying the bag in a different
order leaves the same residue at the bottom --- and that, exactly, is where the granted line, and only
the granted line, does its work.

\section{Real partitions and the finite frame}
\label{se:real-partitions-and-the-finite-frame}

A single spline over a handful of anchors is a toy; real problems live on real domains, and the
construction scales up by partitioning. It divides the domain into a finite set of pieces --- a real
partition --- lays a local spline over each, and assembles the pieces into a finite-element
approximation of the whole. This is the Galerkin method of the seventh chapter, now carried out over
an honest partition of a real interval, and it inherits the same defining virtue: the approximation is
required not to make the residual vanish pointwise, which a finite construction cannot guarantee, but
to make it vanish \emph{weakly} --- orthogonal to every basis function the partition supplies. The
partitioned Lanczos approximation is proved to have exactly this weak-residual-zero property, and that
weak vanishing is the finite-element solution's certificate of correctness.

The whole of this happens inside a single finite frame --- the bounded working memory the seventh
chapter introduced and the construction keeps honest about. The partition is finite, the elements are
finite, the truncation that turns the full problem into the solvable one is a deliberate finite cut,
and the construction never pretends the frame is unbounded. What recovery is possible from such a
finite partition is governed by the sampling law met long ago: the experiment on Fourier and Nyquist
holds here too, that a continuous solution can be reconstructed from the finite partition exactly to
the degree the partition is dense enough to support the inversion. Refine the partition and the
finite-element solution approaches the true one; refine too coarsely and it names something the data
cannot recover. The partition is where the continuum the construction has always only approached is
approached once more, deliberately and with a stated rate.

A concrete partition makes the scaling visible. To solve a problem over a long interval, the
construction does not attempt one enormous spline across the whole of it --- that would overflow the
single finite frame in a few steps. It chops the interval into a finite number of subintervals, fits a
modest local spline on each, and stitches the locals together at the partition boundaries, requiring
only that the pieces agree where they meet and that the assembled residual be weakly zero against the
partition's own basis. Each local problem fits comfortably in the frame; the global solution is their
assembly. This is the finite-element method exactly, and its power is that a problem too large to solve
whole becomes a finite family of problems each small enough to solve, with the partition boundaries
carrying the agreement between them. The cost of the strategy is the cost of the partition: a finer
partition resolves more but spends more frame, and choosing the partition is choosing how much of the
continuum to try to recover.

\section{The interior node and the mixed coupling}
\label{se:the-interior-node-and-the-mixed-coupling}

The last move of the chapter is the one that opens onto the next, and it lives at an interior node. Give
a cubic path a middle node --- an interior joint between a left half and a right half --- and vary it.
The change in the action does not stay tidily on one side. The construction proves that the coupled
difference produced by varying the middle node decomposes into exactly three terms: a residual from the
left half alone, a residual from the right half alone, and a third, the \emph{mixed} term, in which the
left variation and the right variation appear together.

That mixed term is the seed of everything the next chapters are about, and the construction is careful
to carry it whole rather than splitting it away. The first variation, all through Part II, has been a
sum of independent one-sided responses --- vary here, see the change here; vary there, see the change
there. The mixed coupling is the first place two variations are not independent, where moving the left
half changes how the right half responds. That cross-dependence is second-order by nature: it is not
how the action changes, but how one change changes another change, and a quantity built from the
product of two variations is a second variation. The interior node is therefore where the second
variation first appears in the construction, folded into the mixed term of a coupled cubic difference,
and the discipline of keeping that term intact --- never approximating the coupling away as the
first-order analysis happily would --- is what lets the next chapter lift it out and name it. The
single invariant the whole part is named for is sitting in this cross term, waiting.

The decomposition deserves to be seen term by term, because the whole of the next part turns on the
third term. Hold the cubic path's two endpoints fixed and let only the middle node move. The action
changes, and the change splits into three contributions the construction computes separately. The first
is what the left half would have reported on its own, as if the right half were frozen --- the left
middle residual. The second is the mirror image, the right half's own response with the left frozen. If
these were all, the middle node would behave like two independent first variations back to back, and
nothing new would have happened. But there is a third term, alive only because both halves moved at
once: the mixed coupling, in which the left displacement and the right displacement multiply together.
Freeze either half and the mixed term vanishes; it lives only in the simultaneous variation of both. A
quantity proportional to the product of two displacements is, by definition, second-order, and the
construction's refusal to drop it --- to carry the coupled difference with its mixed term intact rather
than approximating it as two independent responses --- is precisely the step that keeps the second
variation in view. The first variation lived in the left and right residuals; the second variation
lives in the mixed one.

\section{Selection by cost}
\label{se:selection-by-cost}

One question remains once the splines are built and the partitions assembled: when several admissible
solutions present themselves, which does the construction take? It selects by cost. Every computation
the construction performs carries a price --- the construction's heartbeat, the running tally of what
each step has spent --- and among admissible solutions it prefers the one reached most cheaply. This is
not an aesthetic preference but a physical one, and the experiment on the thermodynamic cost of erasure
is its anchor: discarding information is not free, it has an irreducible thermodynamic price, so a
solution that throws away less and is reached with fewer erasures is genuinely cheaper in the only
currency that is conserved. Selection by cost is the construction's way of making the variational
principle operational at the level of computation: not merely that the chosen path makes the action
stationary, but that, among the ways of finding it, the construction spends the least.

It is worth saying why cost belongs in a book about measurement at all, and not only in one about
computation. The construction has insisted from its first page that a verdict has a price --- that
deciding whether true equals true costs something, even if almost nothing, and that later decisions
cost more. Selection by cost is that insistence reaching the variational method: of two solutions that
both satisfy the equation, the construction takes the one whose certificate was cheaper to produce, and
it can do this because it has kept the tally all along. The thermodynamic anchor is what makes the
preference more than convention. Erasing a record is not free; it dissipates a definite minimum of
energy, so a computation that erases less is not merely tidier but genuinely costs less in the conserved
currency of the world. The heartbeat the construction listens to is, at bottom, that currency being
spent.

\section{What is demonstrated, and what is assumed}
\label{se:demonstrated-assumed-ch11}

This is the largest chapter, and the separation of built from assumed is, gratifyingly, the same one as
the last chapter --- no new assumption was incurred in all this machinery.

What is demonstrated is substantial. Cubic-spline sufficiency is a proved law: a closed-flat path gives
a cubic spline whose residue is zero. The Lanczos iteration is proved to terminate, by a strictly
decreasing knot count over a finite multiset. The partitioned finite-element approximation is proved to
have a weakly vanishing residual. And the coupled cubic difference is proved to decompose into left,
right, and a mixed term, the second variation's first appearance. Every one of these is a checked
theorem about finite objects.

\begin{quote}
\emph{A note on the one assumption.} Nothing in this chapter adds to the book's single granted step. The
one place a reader might think a new assumption hides --- that the Lanczos iteration reaches the same
terminal residue regardless of the order it consumes knots --- is not a new assumption but the old one
restated: it is the residue-collapse, the universal over directions that the Euler--Lagrange oracle of
the previous chapter was introduced to decide. Termination is built; order-independence is the oracle.
Spline sufficiency, likewise, stands on the equivalence and the oracle of the last chapter, not on
anything new. The honest summary is that this entire chapter of solving is constructive on top of
exactly one assumption, declared once, two chapters ago.
\end{quote}

What is assumed, then, is nothing this chapter introduced. With the equation now solved --- the spline
that suffices, the iteration that computes it, the partition that scales it, and the mixed term that
carries the second variation --- the construction is ready to lift that cross term out of the coupled
difference and look at it directly. The next chapter does exactly that, and the chapter after it gives
the result its name: the second variation is the single invariant, and everything from here to the
field theory is an exploration of what that one quantity controls.
